\begin{figure}[htbp]
\centering
\begin{tikzpicture}[
    font=\small,
    caja/.style={draw=#1, fill=#1!7, rounded corners=2pt, align=center,
                 minimum width=2.9cm, minimum height=0.95cm, inner sep=3pt},
    proceso/.style={caja=subsectionblue},
    exito/.style={draw=subsectionblue, fill=subsectionblue!20, rounded corners=2pt,
                  align=center, minimum width=2.9cm, minimum height=0.95cm,
                  inner sep=3pt, thick},
    rechazo/.style={caja=ugrred},
    aviso/.style={caja=gray},
    decision/.style={chamfered rectangle, chamfered rectangle xsep=4pt,
                     draw=gray!70!black, fill=gray!8, align=center,
                     minimum width=2.9cm, minimum height=0.9cm, inner sep=3pt},
    flecha/.style={-{Stealth[length=2mm]}, thick, draw=gray!70!black},
    et/.style={font=\scriptsize, text=gray!55!black, inner sep=2pt},
]

% --- Camino principal ------------------------------------------------------
\node[proceso] (ini) at (0,0)
      {Petición de Bizum\\[-1pt]\scriptsize regex del backend o \textit{tool call} del LLM};
\node[decision] (d1) at (0,-1.8)  {¿Indica el importe?};
\node[decision] (d2) at (0,-3.6)  {¿Entre 0,50 y 1.000 €?};
\node[decision] (d3) at (0,-5.4)  {¿Existe el contacto?};
\node[proceso]  (pe) at (0,-7.2)
      {Bizum pendiente\\[-1pt]\scriptsize evento \texttt{pedir\_pin}};
\node[decision] (d4) at (0,-9.0)  {¿PIN correcto?};
\node[proceso]  (tx) at (0,-10.8)
      {Transacción \texttt{BEGIN IMMEDIATE}\\[-1pt]\scriptsize límite diario de 500 € y saldo};
\node[exito]    (ok) at (0,-12.6)
      {Descuenta el saldo\\[-1pt]\scriptsize eventos \texttt{saldo} y ticket};

\draw[flecha] (ini) -- (d1);
\draw[flecha] (d1) -- node[et, right] {sí} (d2);
\draw[flecha] (d2) -- node[et, right] {sí} (d3);
\draw[flecha] (d3) -- node[et, right] {sí} (pe);
\draw[flecha] (pe) -- (d4);
\draw[flecha] (d4) -- node[et, right] {sí} (tx);
\draw[flecha] (tx) -- node[et, right] {cabe} (ok);

% --- Salidas a la derecha --------------------------------------------------
\node[aviso]   (r1) at (5.6,-1.8)  {Pregunta el importe\\[-1pt]\scriptsize no deja nada pendiente};
\node[rechazo] (r2) at (5.6,-3.6)  {Rechaza\\[-1pt]\scriptsize sin pedir el PIN};
\node[rechazo] (r3) at (5.6,-5.4)  {Avisa y termina};
\node[aviso]   (r4) at (5.6,-9.0)  {Resta un intento\\[-1pt]\scriptsize 3 como máximo};
\node[rechazo] (c4) at (5.6,-10.8) {Cancela el envío};

\draw[flecha] (d1) -- node[et, above] {no} (r1);
\draw[flecha] (d2) -- node[et, above] {no} (r2);
\draw[flecha] (d3) -- node[et, above] {ninguno} (r3);
\draw[flecha] (d4) -- node[et, above] {no} (r4);
\draw[flecha] (r4.north) |- node[et, right, pos=0.25] {quedan intentos} (pe.east);
\draw[flecha] (r4) -- node[et, right] {sin intentos} (c4);

% --- Salidas a la izquierda ------------------------------------------------
\node[decision] (sg) at (-5.6,-5.4)  {¿Querías decir…?};
\node[rechazo]  (c3) at (-5.6,-3.6)  {Cancela};
\node[rechazo]  (r5) at (-5.6,-10.8) {Rechaza\\[-1pt]\scriptsize la BD no cambia};

\draw[flecha] (d3) -- node[et, above] {parecido} (sg);
\draw[flecha] (sg) -- node[et, right] {no} (c3);
\draw[flecha] (sg.south) |- node[et, right, pos=0.25] {sí} (pe.west);
\draw[flecha] (tx) -- node[et, above] {no cabe} (r5);

% --- Nota sobre el PIN -----------------------------------------------------
\node[font=\scriptsize, text=ugrred, align=center, draw=ugrred!60, dashed,
      rounded corners=2pt, inner sep=4pt] at (-5.6,-8.6)
      {El PIN viaja por el evento\\\texttt{auth\_bizum}: el modelo\\nunca lo recibe};

\end{tikzpicture}
\caption{Flujo de un envío de Bizum. Todas las comprobaciones se hacen en el
backend y en este orden: el importe se valida antes de pedir el PIN, y el
límite diario y el saldo se comprueban dentro de la misma transacción que los
descuenta.}
\label{fig:bizum}
\end{figure}
